\chapter{Receiver and Channel Models}

\section{Continuous-time linear equalizer}

The CTLE is a matched NMOS differential pair with resistive loads, parallel
source-degeneration resistance and capacitance, and an ideal tail-current
source. Its tunable quantities are $R_{\mathrm{LOAD}}$, $R_{\mathrm{DEG}}$,
$C_{\mathrm{DEG}}$, $I_{\mathrm{TAIL}}$, and, in the expanded v3 schema, MOS
width, length, and multiplier. The differential circuit is electrically
inverting, so the benches define the logical output polarity consistently.

At low frequency, $C_{\mathrm{DEG}}$ is approximately open and
$R_{\mathrm{DEG}}$ reduces effective transconductance. At increasing frequency,
the capacitor bypasses part of the degeneration and raises the effective gain.
The zero is approximately
\begin{equation}
  f_z \approx \frac{1}{2\pi R_{\mathrm{DEG}}C_{\mathrm{DEG}}}.
\end{equation}
The complete response also depends on transistor transconductance, output
resistance, load resistance, parasitic capacitance, and the operating point.
Consequently, this expression explains the tuning direction but does not replace
the simulated transfer function.

The current Stage~1 model uses ideal resistors, capacitors, and a tail-current
source. Its computed MOS channel area is only a partial geometry measure; no
total receiver area is claimed.

\section{Four-port channel}

\code{simulator/channel.py} parses Touchstone 1.x four-port data, supports
RI, magnitude-angle, and dB-angle formats, and constructs the differential
transfer $S_{dd21}$. Before transient filtering, the channel is checked for port
mapping, monotonic frequency data, reference impedance, bandwidth, passivity,
and causal behavior. The public development channel and the included synthetic
regression channel are useful test inputs, but neither constitutes PCI-SIG
compliance evidence.

\section{Behavioral one-tap DFE}

For sampled CTLE output $y[n]$, tap $b_1$, and previous decision
$\hat d[n-1]\in\{-1,+1\}$, the DFE uses
\begin{equation}
  u[n] = y[n] - b_1\hat d[n-1], \qquad
  \hat d[n] = \operatorname{sign}(u[n]).
\end{equation}
The implementation first trains or selects a sampling phase, then evaluates
locked-phase eye height, contiguous eye width, decision margin, and observed
errors. Percentile-based vertical metrics reduce sensitivity to single-sample
outliers. The model does not simulate the physical feedback amplifier, slicer,
clock path, or adaptation circuitry.
